\thispagestyle{plain}
\phantomsection
\addcontentsline{toc}{chapter}{Abstract}
\centerline{\zihao{3}\bfseries ABSTRACT}

\linespread{1.4}\zihao{-4}\bigskip

This thesis studies three-dimensional target positioning with multi-station angle and range observations. A unified weighted least-squares model is developed to handle heterogeneous measurements, different noise scales, and angular periodicity. Two numerical solvers are implemented and compared: a Newton-type method using first- and second-order derivatives, and an iterative normal-equation method based on linearization.

To improve convergence robustness, an initial estimate is generated by a three-station geometric trilateration routine before nonlinear optimization. Analytical expressions of residuals, gradients, and Hessians are organized for azimuth, elevation, and slant range observations. A periodic correction strategy is introduced for azimuth residuals to avoid discontinuity near angle wrapping boundaries.

Using the provided course code setting (four angle stations and four range stations), extensive experiments are carried out, including convergence behavior, estimation accuracy, matrix conditioning, uncertainty propagation, confidence ellipsoid construction, and sensitivity to noise and station perturbations. The results show that both methods converge to nearly identical position estimates, while vertical uncertainty is larger than horizontal uncertainty under the current geometry. The inverse normal matrix and the inverse Hessian at convergence present consistent scales and can be used for covariance approximation.

Confidence regions are further constructed from chi-square quantiles and spectral decomposition, providing interpretable uncertainty characterization in principal directions. The thesis also summarizes implementation details and practical parameter recommendations for engineering deployment.

The work delivers an end-to-end pipeline from modeling to derivation, optimization, implementation, and quantitative validation, and can serve as a reusable reference for sensor fusion positioning tasks in radar-electro-optical systems and unmanned platform localization.

\bigskip
\noindent\textbf{\zihao{4} Keywords:}
target positioning; weighted least squares; Newton method; normal equations; confidence ellipsoid; uncertainty analysis
